\section{Определение матрицы Грама}

\begin{definition}[Матрица Грама]
Пусть $V$ --- евклидово ($E$) или унитарное ($U$) пространство, $E = \{e_1, \ldots, e_n\}$ --- базис в $V$.

\textbf{Матрицей Грама} базиса $E$ называется матрица соответствующей билинейной или полуторалинейной формы:
\[
\Gamma = \begin{pmatrix}
(e_1, e_1) & \cdots & (e_1, e_n) \\
\vdots & \ddots & \vdots \\
(e_n, e_1) & \cdots & (e_n, e_n)
\end{pmatrix} = (g_{ij}).
\]
При этом определитель матрицы Грама положителен: $|\Gamma| > 0$.

Для системы векторов $a_1, \ldots, a_k \in V$ ($k \leq n$) \textbf{матрицей типа Грама} называется матрица:
\[
G = \begin{pmatrix}
(a_1, a_1) & \cdots & (a_1, a_k) \\
\vdots & \ddots & \vdots \\
(a_k, a_1) & \cdots & (a_k, a_k)
\end{pmatrix}.
\]
\end{definition}
